\begin{proof}[Proof of \cref{obj:lem:bounded-mean-concentration}]
\begin{enumerate}
\item
Since \(m\ge1\) and a probability space is nonempty, the pointwise bound for \(H_1\) gives
\[
-b\le H_1\le b
\quad\Longrightarrow\quad
-b\le b
\quad\Longrightarrow\quad
0\le b.
\]
If \(b=0\), both asserted bounds follow from
\[
\Pr\left\{m^{-1}\sum_{i=1}^m(H_i-\mathbb EH_i)\ge x\right\}\le1,
\qquad
\Pr\left\{m^{-1}\sum_{i=1}^m(H_i-\mathbb EH_i)\le-x\right\}\le1.
\]
Henceforth suppose \(b>0\), and define the variance proxy
\[
c_b:=\left(\frac{|b-(-b)|}{2}\right)^2=b^2>0.
\]
% lean: CausalSmith.Stat.LdpOptvalueUniformFrontier.boundedMeanConcentration_proved

\item
We first establish an upper-tail bound for any mutually independent, measurable real random variables \(H_1^\star,\ldots,H_m^\star\) on the given probability space satisfying
\[
-b\le H_i^\star\le b,
\qquad i=1,\ldots,m.
\]
Define their centered versions and interval endpoints by
\[
H_i^\circ:=H_i^\star-\mathbb EH_i^\star,
\qquad
\ell_i^\circ:=-b-\mathbb EH_i^\star,
\qquad
u_i^\circ:=b-\mathbb EH_i^\star.
\]
Boundedness ensures integrability, and therefore
\[
\mathbb EH_i^\circ=0,
\qquad
\ell_i^\circ\le H_i^\circ\le u_i^\circ,
\qquad
u_i^\circ-\ell_i^\circ=2b.
\]
The centered variables remain mutually independent because each is a measurable function of its corresponding original variable.

For each \(i\), define the moment-generating function and its logarithm on all of \(\mathbb R\) by
\[
\Phi_i(\upsilon):=\mathbb E\exp(\upsilon H_i^\circ),
\qquad
\psi_i(\upsilon):=\log\Phi_i(\upsilon),
\qquad \upsilon\in\mathbb R.
\]
Boundedness makes \(\Phi_i\) finite and strictly positive and permits differentiation under the expectation. Define expectation under the exponential tilt by
\[
\mathbb E_{\upsilon}^{(i)}[\cdot]
:=
\frac{\mathbb E[\cdot\,\exp(\upsilon H_i^\circ)]}{\Phi_i(\upsilon)},
\qquad
\bar h_{i,\upsilon}:=\mathbb E_{\upsilon}^{(i)}H_i^\circ.
\]
This is expectation under a probability measure absolutely continuous with respect to the original measure. Direct differentiation gives
\[
\psi_i(0)=0,
\qquad
\psi_i'(0)=\mathbb EH_i^\circ=0,
\qquad
\psi_i''(\upsilon)
=
\mathbb E_{\upsilon}^{(i)}[(H_i^\circ)^2]
-\bar h_{i,\upsilon}^{\,2}.
\]

The interval bound persists under the tilt. Consequently,
\[
\begin{aligned}
0
&\le
\mathbb E_{\upsilon}^{(i)}
\bigl[(u_i^\circ-H_i^\circ)(H_i^\circ-\ell_i^\circ)\bigr]\\
&=
-\mathbb E_{\upsilon}^{(i)}[(H_i^\circ)^2]
+(u_i^\circ+\ell_i^\circ)\bar h_{i,\upsilon}
-u_i^\circ\ell_i^\circ.
\end{aligned}
\]
Subtracting the squared tilted mean yields
\[
\begin{aligned}
\psi_i''(\upsilon)
&\le
(u_i^\circ-\bar h_{i,\upsilon})
(\bar h_{i,\upsilon}-\ell_i^\circ)\\
&=
\left(\frac{u_i^\circ-\ell_i^\circ}{2}\right)^2
-
\left(\bar h_{i,\upsilon}
-\frac{u_i^\circ+\ell_i^\circ}{2}\right)^2
\le c_b.
\end{aligned}
\]
For \(\upsilon>0\), Taylor's theorem supplies an intermediate point
\[
\upsilon_{i,\mathrm{mid}}\in(0,\upsilon),
\qquad
\psi_i(\upsilon)
=
\frac{\upsilon^2}{2}\psi_i''(\upsilon_{i,\mathrm{mid}})
\le\frac{c_b\upsilon^2}{2}.
\]
Thus \(\Phi_i(\upsilon)\le\exp(c_b\upsilon^2/2)\) for positive \(\upsilon\); at zero both sides equal one. For negative \(\upsilon\), apply the same positive-parameter calculation to the centered variable \(-H_i^\circ\), whose interval has the same width:
\[
-u_i^\circ\le-H_i^\circ\le-\ell_i^\circ,
\qquad
\mathbb E[-H_i^\circ]=0.
\]
It follows that
\[
\Phi_i(\upsilon)
=
\mathbb E\exp\bigl((-\upsilon)(-H_i^\circ)\bigr)
\le
\exp\left(\frac{c_b(-\upsilon)^2}{2}\right).
\]
We have therefore established
\[
\mathbb E\exp(\upsilon H_i^\circ)
\le\exp\left(\frac{c_b\upsilon^2}{2}\right),
\qquad \upsilon\in\mathbb R,\quad i=1,\ldots,m.
\]
% lean: ProbabilityTheory.hasSubgaussianMGF_of_mem_Icc

\item
Define the centered sum
\[
S^\star:=\sum_{i=1}^m H_i^\circ.
\]
Independence gives the factorization
\[
\begin{aligned}
\mathbb E\exp(\upsilon S^\star)
&=\prod_{i=1}^m\mathbb E\exp(\upsilon H_i^\circ)\\
&\le
\prod_{i=1}^m\exp\left(\frac{c_b\upsilon^2}{2}\right)
=
\exp\left(\frac{mc_b\upsilon^2}{2}\right),
\qquad \upsilon\in\mathbb R.
\end{aligned}
\]
Indeed, adjoining one centered variable to a partial sum multiplies their moment-generating functions, since that variable is independent of the preceding sum; induction gives the displayed formula.

Since \(mx\ge0\), for every \(\upsilon\ge0\) exponential Markov's inequality gives
\[
\begin{aligned}
\Pr\{S^\star\ge mx\}
&\le
\exp(-\upsilon mx)\,\mathbb E\exp(\upsilon S^\star)\\
&\le
\exp\left(-\upsilon mx+\frac{mc_b\upsilon^2}{2}\right).
\end{aligned}
\]
Choose
\[
\upsilon_{\mathrm{opt}}:=\frac{mx}{mc_b}=\frac{x}{c_b}\ge0.
\]
This choice is defined also when \(x=0\), and substitution yields
\[
\Pr\{S^\star\ge mx\}
\le
\exp\left(-\frac{(mx)^2}{2mc_b}\right).
\]
% lean: ProbabilityTheory.HasSubgaussianMGF.measure_sum_ge_le_of_iIndepFun

\item
Because \(m>0\), multiplying by \(m\) gives the event identity
\[
\left\{
x\le m^{-1}\sum_{i=1}^m
(H_i^\star-\mathbb EH_i^\star)
\right\}
=
\{mx\le S^\star\}.
\]
Using \(c_b=b^2\), the exponent simplifies to
\[
-\frac{(mx)^2}{2mc_b}
=
-\frac{mx^2}{2b^2}.
\]
Hence the upper-tail calculation establishes
\[
\Pr\left\{
m^{-1}\sum_{i=1}^m(H_i^\star-\mathbb EH_i^\star)\ge x
\right\}
\le
\exp\left(-\frac{mx^2}{2b^2}\right).
\]
Taking \(H_i^\star=H_i\) proves the asserted upper-tail bound.
% lean: CausalSmith.Stat.LdpOptvalueUniformFrontier.boundedMeanConcentration_proved.upper

\item
For the lower tail, define
\[
H_i^-:=-H_i,
\qquad i=1,\ldots,m.
\]
These are mutually independent and measurable, and negating the original bounds gives
\[
-b\le H_i^-\le b.
\]
Linearity of expectation and finite summation gives
\[
\mathbb EH_i^-=-\mathbb EH_i,
\qquad
\sum_{i=1}^m(H_i^--\mathbb EH_i^-)
=
-\sum_{i=1}^m(H_i-\mathbb EH_i).
\]
Therefore,
\[
\left\{
m^{-1}\sum_{i=1}^m(H_i-\mathbb EH_i)\le-x
\right\}
=
\left\{
m^{-1}\sum_{i=1}^m(H_i^--\mathbb EH_i^-)\ge x
\right\}.
\]
Applying the upper-tail bound just established to \(H_1^-,\ldots,H_m^-\) yields
\[
\Pr\left\{
m^{-1}\sum_{i=1}^m(H_i-\mathbb EH_i)\le-x
\right\}
\le
\exp\left(-\frac{mx^2}{2b^2}\right),
\]
which completes the proof.
% lean: CausalSmith.Stat.LdpOptvalueUniformFrontier.boundedMeanConcentration_proved
\end{enumerate}
\end{proof}
